\documentclass[10pt,a4paper]{amsart}
\usepackage[margin=19mm]{geometry}
\usepackage{amsmath,amssymb,mathtools}
\usepackage[scr=rsfs,cal=euler]{mathalfa}
\usepackage{xfrac}
\usepackage[unicode,hidelinks,bookmarks=false]{hyperref}
\newcommand{\ie}{i.e.\ }
\newcommand{\cf}{cf.\ }
\newcommand{\resp}{respectively\ }
\newcommand{\Subsection}{\subsection}
\providecommand{\nobreakdash}{\nobreak\hbox{-}}
\setlength{\emergencystretch}{2em}
\multlinegap0pt
\usepackage{bm}
\usepackage{float}
\usepackage{amssymb}
\newtheorem{prop}{}
\title{On the sign changes of $\psi(x)-x$}
\author{Maciej Grze\'skowiak, Jerzy Kaczorowski, {\L}ukasz Pa\'nkowski, Maciej Radziejewski}
\date{Source: arXiv:2408.10399v2 (CC BY 4.0)}
\begin{document}
\maketitle

\noindent\emph{Source and licence.} On the sign changes of $\psi(x)-x$. Version arXiv:2408.10399v2 (CC BY 4.0), distributed by arXiv under the Creative Commons Attribution 4.0 International licence, http://creativecommons.org/licenses/by/4.0/, as recorded in the arXiv licence metadata for this exact version and shown on the abstract page https://arxiv.org/abs/2408.10399v2. Full licence notice: \texttt{licenses/arxiv-2408.10399-CC-BY-4.0.txt}.\par\smallskip

\noindent\textit{Editorial scope.} Counted unit: the article's prop prop:sufficient-condition (source0.tex lines 197--234). The article's complete original proof of it is delivered with it (source0.tex lines 360--410, one \texttt{proof} environment of 51 lines, which the article places away from the statement). No mathematics is written, repaired or re-derived: the statement and the proof are the article's own lines, in the article's own notation, and no retained line is substituted or re-flowed.  No further source-local context is needed: every cross-reference of every retained line resolves inside the two delivered spans.  Nothing was omitted from any retained span, so the package carries no omission record.  No retained line contains a citation, so the wrapper carries no bibliography and no bibliography entry.  No retained line contains a figure environment, an image-inclusion command or drawing code, so no image is delivered and the package declares no asset dependency.  Editorial formatting only, in addition: an AMS \texttt{amsart} preamble that restates the article's own declarations and macros used by the retained lines, namely the environments \texttt{prop} on separate counters, together with the standard packages those lines need.  The retained environments print in this document as Proposition 1 in order of appearance, whereas the article prints the same items with the numbering of its own full document.  The complete native source of the article as distributed in the arXiv e-print for this exact version is bundled unchanged as \texttt{sources/164001/source0.tex}.
\begin{prop}
\label{prop:sufficient-condition}Let $\boldsymbol{\eta}=(\eta_{1},\dotsc,\eta_{N})\in\mathbb{R}^{N}$
and denote
\[
F_{N}^{\boldsymbol{\eta}}(z)=1+\sum_{n=1}^{N}a(n)e^{i\left(\eta_{n}+\omega_{n}z\right)}.
\]
Suppose $z_{1}(t)$, $0\leq t\leq1$, is a closed contour in $\mathbb{H}$,
oriented counterclockwise, not passing through any zero of $F_{N}^{\boldsymbol{\eta}}(z)$.
Suppose $z_{2}(t)$ is a closed contour around $0$, also oriented
counterclockwise, satisfying $\left|z_{2}(t)\right|=1$, $0\leq t\leq1$,
and $\eta:[0,1]\to[-\tfrac{\pi}{2},\tfrac{\pi}{2}]$ is an arbitrary
function. Let
\[
\varphi_{n}(t)=\eta_{n}+\arg a(n)+\omega_{n}\Re z_{1}(t)-\arg z_{2}(t)-\eta(t),\qquad0\leq t\leq1,n=1,\dotsc,N.
\]
Suppose $u(t)$ is a function satisfying
\[
0<u(t)<\Re\left(F_{N}^{\boldsymbol{\eta}}(z_{1}(t))\overline{z_{2}(t)}e^{-i\eta(t)}\right)-R_{N}(\Im z_{1}(t)),\qquad0\leq t\leq1.
\]
For $\varphi,\psi\in[0,\pi]$ let 
\[
v(\varphi,\psi)=\begin{cases}
\cos\left(\varphi\right)-\cos\left(\varphi+\psi\right), & \varphi+\psi\leq\pi,\\
\cos\left(\varphi\right)+1, & \text{otherwise.}
\end{cases}
\]
Suppose $w_{n}:[0,\tfrac{1}{2}]\to[0,1]$ are non-decreasing functions
such that
\begin{equation}
w_{n}(x)\geq\sup_{0\leq t\leq1}\frac{\left|a(n)\right|}{u(t)e^{\omega_{n}\Im z_{1}(t)}}v\left(2\pi\left\Vert \frac{\varphi_{n}(t)}{2\pi}\right\Vert ,2\pi x\right)\label{eq:wn(x)-assumption}
\end{equation}
for every $x\in[0,\tfrac{1}{2}]$ and $n=1,\dotsc,N$. Then for every
$\tau\in\mathbb{R}$ such that
\begin{equation}
\sum_{n=1}^{N}w_{n}\left(\left\Vert \frac{\omega_{n}\tau-\eta_{n}}{2\pi}\right\Vert \right)\leq1\label{eq:prop-sum-of-penalties}
\end{equation}
there is a zero of $F(z)$ inside the contour $z_{1}(t)+\tau$.
\end{prop}

\begin{proof}[Proof of Proposition~\ref{prop:sufficient-condition}]
Fix $\tau\in\mathbb{R}$ that satisfies (\ref{eq:prop-sum-of-penalties}).
By the argument principle it is enough to show that 
\[
F(z_{1}(t)+\tau)\neq-z_{2}(t)\left|F(z_{1}(t)+\tau)\right|,\qquad0\leq t\leq1,
\]
which follows from
\[
\Re\left(F(z_{1}(t)+\tau)\overline{z_{2}(t)}e^{-i\eta(t)}\right)>0,\qquad0\leq t\leq1.
\]
We have
\begin{align*}
\Re\left(F(z_{1}(t)+\tau)\overline{z_{2}(t)}e^{-i\eta(t)}\right) & \geq\Re\left(F_{N}(z_{1}(t)+\tau)\overline{z_{2}(t)}e^{-i\eta(t)}\right)-R_{N}(\Im z_{1}(t))\\
 & >\Re\left(\Bigl(F_{N}(z_{1}(t)+\tau)-F_{N}^{\boldsymbol{\eta}}(z_{1}(t))\Bigr)\overline{z_{2}(t)}e^{-i\eta(t)}\right)+u(t),
\end{align*}
so it is enough to show
\begin{equation}
\Re\left(\Bigl(-F_{N}(z_{1}(t)+\tau)+F_{N}^{\boldsymbol{\eta}}(z_{1}(t))\Bigr)\overline{z_{2}(t)}e^{-i\eta(t)}\right)\leq u(t),\qquad0\leq t\leq1.\label{eq:sufficient-condition-u(t)}
\end{equation}
We have

\begin{align*}
\Re\biggl(\Bigl(-F_{N}(z_{1}(t)+\tau) & +F_{N}^{\boldsymbol{\eta}}(z_{1}(t))\Bigr)\overline{z_{2}(t)}e^{-i\eta(t)}\biggr)\\
= & \Re\sum_{n=1}^{N}a(n)e^{i\left(\eta_{n}+\omega_{n}z_{1}(t)-\arg z_{2}(t)-\eta(t)\right)}\left(1-e^{i\left(\omega_{n}\tau-\eta_{n}\right)}\right)\\
= & \sum_{n=1}^{N}\left|a(n)\right|e^{-\omega_{n}\Im z_{1}(t)}\left(\cos\left(\varphi_{n}(t)\right)-\cos\left(\omega_{n}\tau-\eta_{n}+\varphi_{n}(t)\right)\right).
\end{align*}
Since 
\[
\cos\left(\omega_{n}\tau-\eta_{n}+\varphi_{n}(t)\right)=\cos\left(2\pi\left(\left\Vert \frac{\omega_{n}\tau-\eta_{n}+\varphi_{n}(t)}{2\pi}\right\Vert \right)\right)
\]
and
\[
\left\Vert \frac{\omega_{n}\tau-\eta_{n}+\varphi_{n}(t)}{2\pi}\right\Vert \leq\left\Vert \frac{\omega_{n}\tau-\eta_{n}}{2\pi}\right\Vert +\left\Vert \frac{\varphi_{n}(t)}{2\pi}\right\Vert ,
\]
moreover $\cos(x)$ is monotone decreasing on $[0,\pi]$, we have
\[
\cos\left(\omega_{n}\tau-\eta_{n}+\varphi_{n}(t)\right)\geq\cos\left(2\pi\left(\left\Vert \frac{\omega_{n}\tau-\eta_{n}}{2\pi}\right\Vert +\left\Vert \frac{\varphi_{n}(t)}{2\pi}\right\Vert \right)\right)
\]
when $\left\Vert \frac{\omega_{n}\tau-\eta_{n}}{2\pi}\right\Vert +\left\Vert \frac{\varphi_{n}(t)}{2\pi}\right\Vert \leq\frac{1}{2}$
and $\cos\left(\omega_{n}\tau-\eta_{n}+\varphi_{n}(t)\right)\geq-1$
in any case. Hence
\[
\cos\left(\varphi_{n}(t)\right)-\cos\left(\omega_{n}\tau-\eta_{n}+\varphi_{n}(t)\right)\leq v\left(2\pi\left\Vert \frac{\varphi_{n}(t)}{2\pi}\right\Vert ,2\pi\left\Vert \frac{\omega_{n}\tau-\eta_{n}}{2\pi}\right\Vert \right)
\]
and
\begin{align*}
\Re\biggl(\Bigl(-F_{N}(z_{1}(t)+\tau) & +F_{N}^{\boldsymbol{\eta}}(z_{1}(t))\Bigr)\overline{z_{2}(t)}\biggr)\\
\leq & \sum_{n=1}^{N}\left|a(n)\right|e^{-\omega_{n}\Im z_{1}(t)}v\left(2\pi\left\Vert \frac{\varphi_{n}(t)}{2\pi}\right\Vert ,2\pi\left\Vert \frac{\omega_{n}\tau-\eta_{n}}{2\pi}\right\Vert \right).
\end{align*}
This implies~(\ref{eq:sufficient-condition-u(t)}) and the assertion. 
\end{proof}

\end{document}
